\section{Threats to Validity}
\label{sec:threats}

Each threat is reported in the same four parts: what it is, how large it is, the diagnostic or
mitigation that answers it, and the residual.

\subsection{Construct validity: what pass-rate measures}

\textbf{Threat.} Pass-rate is a lower bound on semantic equivalence, and it has a floor: a
candidate written without consulting the decompiled program can still pass some tests, so part
of any score may be pattern-matching rather than reconstruction.

\textbf{Quantification.} The floor is measurable inside the study rather than assumed, and
\S\ref{sec:results} reports it: the runs that called no tool at all average 0.113 against 0.638 for those
that called at least one.

\textbf{Residual.} The floor exists and does not dominate, and both halves carry a number. Pass-rate is
reported as a lower bound throughout, and no claim in this paper requires it to be more.

\subsection{Construct validity: what a transport carries}

\textbf{Threat.} The transport contrast is a change of \emph{format} and also a change in the volume of
context the request carries. \S\ref{sec:runtime} reports the same quantity between two services; it also
differs between the two transports, on the axis this study exists to isolate.

\textbf{Quantification.} Median provider-reported input tokens per run, native over textual, across the
eight cells where both arms exist: $2.89$ and $2.47$ on gpt-oss (Databricks, Bedrock), $2.34$ and $1.02$ on
llama, $1.38$ and $1.37$ on haiku, and $0.75$ and $0.73$ on sonnet. The native arm is reported higher in six
cells of eight and lower in two, so the divergence has no fixed direction. T1 is the cell where this matters
most: gpt-oss issues exactly one call per turn, so it is the contrast with no batching confound to absorb
the difference.

\textbf{Diagnostic.} This is a component of the transport as evaluated, not a leak into it. Under native
calling the tool schemas travel in a \texttt{tools} field and results return in the provider's tool role;
under our prompted-text protocol both travel as ordinary messages. \textbf{The two implementations we
evaluated do not isolate serialisation format from context volume}, and the ablation arm of
\S\ref{sec:rq5} separates batching from format but not format from volume. Whether some other pairing of
implementations could hold context volume fixed while changing the encoding is a question this design does
not answer and does not foreclose.

\textbf{Residual.} The transport effects of \S\ref{sec:results} are therefore effects of a transport as
deployed rather than of a serialisation syntax holding all else equal, and a reader sizing their own study
should read them that way.

\subsection{Construct validity: what the ablation arm separates}

\textbf{Threat.} The constrained-native arm of \S\ref{sec:rq5} separates call batching from call format
only up to the two effects below.

\textbf{Quantification.} The constraint does not remove batching alone: forcing one call per turn also
changes the shape of the trajectory the model plans, so what the arm bounds is an upper limit on the
batching contribution rather than an estimate of it. And it covers one infrastructure and the two models
that batch at all --- on the other two the rate is exactly 1.000 calls per turn, so there is nothing to
remove and no contrast to run.

\textbf{Residual.} The arm therefore supports a ceiling rather than a point estimate --- at most 26\% and
13\% of the transport effect on the two models that batch (\S\ref{sec:rq5}) --- and the artifact carries
both bounds with the per-model rates.

\subsection{Internal validity: a working directory shared across cells}

\textbf{Threat.} Candidates were compiled in a directory keyed on the program and the run index and
nothing else, so under concurrent collection two cells could share a compilation path and a run compiling
another cell's candidate would be scored normally with no flag raised. The key omits the cell, so no
contrast is structurally exempt.

\textbf{Quantification.} The precondition --- two cells sharing a compilation path at the same moment ---
arose \textbf{thirty} times in the original batch of 5{,}805 measurements. The deposit carries the
mechanism, the write-to-read window and the probe these data admit.

\textbf{Mitigation.} The working directory now carries the cell, and the whole design was collected again
on workspaces isolated per cell. \textbf{That is what makes the re-collection the primary basis: the
mechanism is absent there by construction, not observed absent} --- thirty occurrences originally, zero
under 451 concurrent measurements that would have exposed it. This is a property of the code rather than
of either batch's numbers.

\textbf{Concordance with the original collection.} Four diagnostics were specified for the comparison,
and all four are met: six of eight contrasts keep their sign against a threshold of six; no test passes
its Holm threshold in either batch; seven of eight re-collected effects fall inside the original 95\%
interval, the exception being T7, whose original interval spans 0.3pp because that cell is almost
noiseless; and the noise share that inverted the sample-size calibration reproduces within an order of
magnitude, 0.1\%/$K_\infty{=}705$ against 0.5\%/668. Appendix~\ref{sec:appendice-provenienza} states when
each threshold was written and against what was known.

% Generata da analysis/tabella_riraccolta.py — non modificare a mano.
\begin{table}[t]
\centering
\caption{The same eight contrasts measured twice: the original collection, and an independent re-collection on workspaces isolated per cell. Differences in percentage points; the last column states whether the re-collected effect falls inside the original 95\% interval. The criteria and their thresholds were written before the comparison could be computed; their timing is stated precisely in \S\ref{sec:threats}, where it is weaker than the pre-registration's.}
\label{tab:riraccolta}
\footnotesize
\setlength{\tabcolsep}{5pt}
\begin{tabular}{@{}llrrc@{}}
\toprule
 & model & orig. & re-coll. & in CI \\
\midrule
T1 & gpt-oss-120b & $-6.4$ & $-6.6$ & yes \\
T2 & llama-3.3-70b & $+3.5$ & $+2.2$ & yes \\
T3 & claude-haiku-4-5 & $-10.4$ & $-9.8$ & yes \\
T4 & claude-sonnet-4-5 & $+2.1$ & $-1.0$ & yes \\
T5 & gpt-oss-120b & $+3.4$ & $+3.3$ & yes \\
T6 & llama-3.3-70b & $-8.9$ & $-7.9$ & yes \\
T7 & claude-haiku-4-5 & $-0.1$ & $+0.2$ & no \\
T8 & claude-sonnet-4-5 & $-0.7$ & $-0.9$ & yes \\
\bottomrule
\end{tabular}
\end{table}

% criterio 1, segni concordi: 6/8 (soglia 6) -> invariato
% criterio 3, copertura IC95: 7/8 (soglia 6) -> invariato


\textbf{Residual.} The original batch is reported in full beside the primary one rather than invalidated:
the criterion this programme applies to invalidation is a systematic, arm-asymmetric channel
(\S\ref{sec:discussion}), and the shared workspace is neither.

\subsection{Conclusion validity: the guards, and their perimeter}

\textbf{Threat and residual.} The analysis is protected by a content-hash chain over six named files ---
the pre-registration, the frozen analysis and power scripts, the held-out corpus, the collector, and the
comparison script --- and by an append-only convention on the result directory whose guarantee is the
committed history. Both perimeters are stated so that a claim of precedence can be checked against the
right one, and the artifact enumerates them.

\subsection{Conclusion validity: the estimator}

\textbf{Threat.} The pre-registered test is a paired $t$ over per-binary means, which discards the
within-binary structure a run-level model would use.

\textbf{Quantification.} What we can measure is where the variance sits, and Table~\ref{tab:variance}
measures it: on T3 and T6, the two contrasts on which the variance decomposition is most informative,
run-to-run noise is 0.5\% and 9.1\% of the
paired variance, so almost all of it is between binaries rather than within them. On T5 it is 93\%, and on
T7 and T8 the estimator does not fit at all.

\textbf{Diagnostic.} The run-level counterfactual is computed in Appendix~\ref{sec:appendice-glmm} and
agrees: a marginal model with cluster-robust errors reproduces the frozen test's reading on all three
contrasts it covers, and none passes Holm.

\textbf{Residual.} That model is reported as a sensitivity analysis rather than substituted for the frozen
estimator: choosing an estimator after seeing which contrasts it favours is the freedom the freeze exists
to remove.

\subsection{External validity: what the design does not settle}

\textbf{Threat.} The design manipulates transport and infrastructure at a fixed \emph{nominal model
identity}, and the weaker phrase is deliberate: the design does not verify that a nominally identical
model is byte-identical across two providers, and neither provider exposes enough to settle it.

\textbf{Quantification.} The divergence matters rather than being hypothetical: on the one axis where the
Wilson bands separate, the same nominal model returns identical scores on 24.4\% of binaries at one cloud
against 66.7\% at the other (\S\ref{sec:results}).

\textbf{Residual.} The rates are reported at the level the design identifies --- the deployed service ---
since attributing them further would require knowing what each provider serves. The scope of every effect
size reported here is one task family, four models and two clouds, as the pre-registration bounds it.
